% GENERATED FILE. Edit canonical lessons, then run npm run generate:books.
% canonical-source-hash: fnv1a64:30fcacd72408ce43
% canonical-lessons: SA-C311-kimvadanti, SA-C311-atankah, SA-C311-uttejayati, SA-C311-sankate, SA-C311-bhrantih

\chapter{A Rumour Spreads}
\label{ch:sa-a-rumour-spreads}
\hlchaptermodality{311}

\begin{chapteropening}
\textbf{By the end of this chapter:} \emph{I can say how a rumour spread and how people reacted.}

Complete the last lesson of chapter 311: I can say how a rumour spread and how people reacted.
\end{chapteropening}

\section[\texorpdfstring{\sk{किंवदन्ती} (kiṃvadantī)}{किंवदन्ती (kiṃvadantī)}]{\sk{किंवदन्ती} (kiṃvadantī) — a rumour}

\label{lesson:SA-C311-kimvadanti}



\begin{quote}
Before the new one: say the Sanskrit for publishes, then the Sanskrit for a detail.
\end{quote}

\subsection*{You'll want to know: \sk{किंवदन्ती}}
\textbf{\sk{किंवदन्ती}} — \emph{kiṃvadantī} — \textquotedblleft{}a rumour\textquotedblright{}.

\begin{quote}
\sk{नगरे} \sk{एका} \sk{किंवदन्ती} \sk{प्रसरति।}
\end{quote}

\begin{quote}
\emph{nagare ekā kiṃvadantī prasarati.}
\end{quote}

\begin{quote}
A rumour is spreading in the city.
\end{quote}

A feminine noun.

\begin{cousinweb}
\textbf{\sk{किंवदन्ती}} (\emph{kiṃvadantī}) is built from \emph{kim}, \textquotedblleft{}what\textquotedblright{}, and a form of the root √vad, \textquotedblleft{}to say\textquotedblright{}: MW explains it as the question what do they say, common talk, rumour (MW).
\end{cousinweb}

\begin{grammarlens}[title={What you've built}]
One more word for telling what happened.
\end{grammarlens}

\subsection*{Your turn}
\begin{itemize}
\raggedright
  \item \emph{Say it:} \emph{kiṃvadantī}
  \item \emph{Say it:} \emph{kiṃvadantī}, once more
  \item \emph{Say it:} \emph{nagare ekā kiṃvadantī prasarati.} — A rumour is spreading in the city.
  \item \emph{Recall:} say the Sanskrit for publishes, then the Sanskrit for a detail, then say \emph{kiṃvadantī} again
\end{itemize}

\begin{tcolorbox}[breakable,colback=teal!4,colframe=teal!35!black,title={Before you move on}]
What does \sk{किंवदन्ती} mean? (\textquotedblleft{}A rumour\textquotedblright{}.) Say it once more.
\end{tcolorbox}
\section[\texorpdfstring{\sk{आतंकः} (ātaṅkaḥ)}{आतंकः (ātaṅkaḥ)}]{\sk{आतंकः} (ātaṅkaḥ) — dread; panic}

\label{lesson:SA-C311-atankah}



\begin{quote}
Before the new one: say the Sanskrit for a detail, then the Sanskrit for a rumour.
\end{quote}

\subsection*{You'll want to know: \sk{आतंकः}}
\textbf{\sk{आतंकः}} — \emph{ātaṅkaḥ} — \textquotedblleft{}dread; panic\textquotedblright{}.

\begin{quote}
\sk{किंवदन्त्या} \sk{जनेषु} \sk{आतंकः} \sk{वर्धते।}
\end{quote}

\begin{quote}
\emph{kiṃvadantyā janeṣu ātaṅkaḥ vardhate.}
\end{quote}

\begin{quote}
Because of the rumour dread grows among the people.
\end{quote}

A masculine noun.

\begin{cousinweb}
\textbf{\sk{आतंकः}} (\emph{ātaṅkaḥ}) MW gives \emph{ātaṅka} as disease, and also as affliction of mind, apprehension, fear, without a root in the entry, so its origin is left uncertain here (MW).
\end{cousinweb}

\begin{grammarlens}[title={What you've built}]
One more word for telling what happened.
\end{grammarlens}

\subsection*{Your turn}
\begin{itemize}
\raggedright
  \item \emph{Say it:} \emph{ātaṅkaḥ}
  \item \emph{Say it:} \emph{ātaṅkaḥ}, once more
  \item \emph{Say it:} \emph{kiṃvadantyā janeṣu ātaṅkaḥ vardhate.} — Because of the rumour dread grows among the people.
  \item \emph{Recall:} say the Sanskrit for a detail, then the Sanskrit for a rumour, then say \emph{ātaṅkaḥ} again
\end{itemize}

\begin{tcolorbox}[breakable,colback=teal!4,colframe=teal!35!black,title={Before you move on}]
What does \sk{आतंकः} mean? (\textquotedblleft{}Dread; panic\textquotedblright{}.) Say it once more.
\end{tcolorbox}
\section[\texorpdfstring{\sk{उत्तेजयति} (uttejayati)}{उत्तेजयति (uttejayati)}]{\sk{उत्तेजयति} (uttejayati) — stirs up; excites}

\label{lesson:SA-C311-uttejayati}



\begin{quote}
Before the new one: say the Sanskrit for a rumour, then the Sanskrit for dread.
\end{quote}

\subsection*{You'll want to know: \sk{उत्तेजयति}}
\textbf{\sk{उत्तेजयति}} — \emph{uttejayati} — \textquotedblleft{}stirs up; excites\textquotedblright{}.

\begin{quote}
\sk{धूर्तः} \sk{जनान्} \sk{उत्तेजयति।}
\end{quote}

\begin{quote}
\emph{dhūrtaḥ janān uttejayati.}
\end{quote}

\begin{quote}
The rogue stirs up the people.
\end{quote}

Past (he …): \sk{सः} \sk{जनान्} \sk{उत्तेजितवान्}, \emph{saḥ janān uttejitavān}.

\begin{cousinweb}
\textbf{\sk{उत्तेजयति}} (\emph{uttejayati}) is built from \emph{ud}, \textquotedblleft{}up\textquotedblright{}, and the causative of the root √tij, \textquotedblleft{}to be sharp\textquotedblright{}; MW gives \emph{ut-tejayati} as to excite, stimulate, incite (MW).
\end{cousinweb}

\begin{grammarlens}[title={What you've built}]
One more word for telling what happened.
\end{grammarlens}

\subsection*{Your turn}
\begin{itemize}
\raggedright
  \item \emph{Say it:} \emph{uttejayati}
  \item \emph{Say it:} \emph{uttejayati}, once more
  \item \emph{Say it:} \emph{dhūrtaḥ janān uttejayati.} — The rogue stirs up the people.
  \item \emph{Recall:} say the Sanskrit for a rumour, then the Sanskrit for dread, then say \emph{uttejayati} again
\end{itemize}

\begin{tcolorbox}[breakable,colback=teal!4,colframe=teal!35!black,title={Before you move on}]
What does \sk{उत्तेजयति} mean? (\textquotedblleft{}Stirs up; excites\textquotedblright{}.) Say it once more.
\end{tcolorbox}
\section[\texorpdfstring{\sk{शंकते} (śaṅkate)}{शंकते (śaṅkate)}]{\sk{शंकते} (śaṅkate) — doubts; suspects}

\label{lesson:SA-C311-sankate}



\begin{quote}
Before the new one: say the Sanskrit for dread, then the Sanskrit for stirs up.
\end{quote}

\subsection*{You'll want to know: \sk{शंकते}}
\textbf{\sk{शंकते}} — \emph{śaṅkate} — \textquotedblleft{}doubts; suspects\textquotedblright{}.

\begin{quote}
\sk{मीरा} \sk{किंवदन्तीं} \sk{शंकते।}
\end{quote}

\begin{quote}
\emph{mīrā kiṃvadantīṃ śaṅkate.}
\end{quote}

\begin{quote}
Mira doubts the rumour.
\end{quote}

Past (she …): \sk{सा} \sk{किंवदन्तीं} \sk{शंकितवती}, \emph{sā kiṃvadantīṃ śaṅkitavatī}.

\begin{cousinweb}
\textbf{\sk{शंकते}} (\emph{śaṅkate}) comes from the root √śaṅk, \textquotedblleft{}to be anxious about, doubt, distrust\textquotedblright{} (MW).
\end{cousinweb}

\begin{grammarlens}[title={What you've built}]
One more word for telling what happened.
\end{grammarlens}

\subsection*{Your turn}
\begin{itemize}
\raggedright
  \item \emph{Say it:} \emph{śaṅkate}
  \item \emph{Say it:} \emph{śaṅkate}, once more
  \item \emph{Say it:} \emph{mīrā kiṃvadantīṃ śaṅkate.} — Mira doubts the rumour.
  \item \emph{Recall:} say the Sanskrit for dread, then the Sanskrit for stirs up, then say \emph{śaṅkate} again
\end{itemize}

\begin{tcolorbox}[breakable,colback=teal!4,colframe=teal!35!black,title={Before you move on}]
What does \sk{शंकते} mean? (\textquotedblleft{}Doubts; suspects\textquotedblright{}.) Say it once more.
\end{tcolorbox}
\section[\texorpdfstring{\sk{भ्रान्तिः} (bhrāntiḥ)}{भ्रान्तिः (bhrāntiḥ)}]{\sk{भ्रान्तिः} (bhrāntiḥ) — a misunderstanding; a false belief}

\label{lesson:SA-C311-bhrantih}



\begin{quote}
Before the new one: say the Sanskrit for stirs up, then the Sanskrit for doubts.
\end{quote}

\subsection*{You'll want to know: \sk{भ्रान्तिः}}
\textbf{\sk{भ्रान्तिः}} — \emph{bhrāntiḥ} — \textquotedblleft{}a misunderstanding; a false belief\textquotedblright{}.

\begin{quote}
\sk{सा} \sk{किंवदन्ती} \sk{केवलं} \sk{भ्रान्तिः} \sk{अस्ति।}
\end{quote}

\begin{quote}
\emph{sā kiṃvadantī kevalaṃ bhrāntiḥ asti.}
\end{quote}

\begin{quote}
That rumour is only a misunderstanding.
\end{quote}

A feminine noun.

\begin{cousinweb}
\textbf{\sk{भ्रान्तिः}} (\emph{bhrāntiḥ}) comes from the root √bhram, \textquotedblleft{}to wander, roam\textquotedblright{}, as in \emph{bhramati}; MW gives \emph{bhrānti} as wandering, and also as perplexity, error, false opinion (MW).
\end{cousinweb}

\begin{grammarlens}[title={What you've built}]
One more word for telling what happened.
\end{grammarlens}

\subsection*{Your turn}
\begin{itemize}
\raggedright
  \item \emph{Say it:} \emph{bhrāntiḥ}
  \item \emph{Say it:} \emph{bhrāntiḥ}, once more
  \item \emph{Say it:} \emph{sā kiṃvadantī kevalaṃ bhrāntiḥ asti.} — That rumour is only a misunderstanding.
  \item \emph{Recall:} say the Sanskrit for stirs up, then the Sanskrit for doubts, then say \emph{bhrāntiḥ} again
\end{itemize}

\begin{tcolorbox}[breakable,colback=teal!4,colframe=teal!35!black,title={Before you move on}]
What does \sk{भ्रान्तिः} mean? (\textquotedblleft{}A misunderstanding; a false belief\textquotedblright{}.) Say it once more.
\end{tcolorbox}
